\section{Related Work}\label{sec:related}

\paragraph{The library.} The developments use the Agda library of our
earlier paper~\cite{anonymous2027popl}. That paper is about the library
itself: presentations as congruences on words, normal forms as bundles
of the standard library, coset tables and towers in the style of
Reidemeister and Schreier, products and extensions of presentations,
and circuits indexed by their width. It tested the design on symmetric
groups, wreath products, Pauli groups, and the single-qubit and
single-qutrit Clifford+$T$ monoids. It also compares the library with
work on presentations and rewriting in Lean, Rocq and
Isabelle~\cite{mathlib-doc-presented,gonthier2013oddorder,kharim2023freegroups,sternagel2013kb},
and we do not repeat that comparison here. The present paper is about
using the library on published proofs that were not written with
formalization in mind, and about what had to be added to it for that:
the matrix layer and the independence framework of
Section~\ref{sec:scaffold}, which take the library across the line its
own limitations section drew.

\paragraph{Bian and Selinger's Agda proofs.} Bian and Selinger checked
part of their two- and three-qubit results in
Agda~\cite{bian2023cliffordt2,bian2023cliffordcs3,bian2023thesis}:
that the circuit relations and a set of two-level relations present
the same group, along a translation of generators, by the
Reidemeister--Schreier method, with generated proof terms. The other
part, that the two-level relations present a matrix group, is
Greylyn's thesis~\cite{greylyn2014} and Bian and Selinger's paper on
$U_n(\Z[\frac12,i])$~\cite{bian2021un}, and was proved on paper. We
formalized both two-level presentations and ported their code, so that
both results now end in a statement about matrices.

\paragraph{Verified quantum programming.} Several projects verify
quantum programs and compilers against a matrix semantics: QWIRE and
SQIR, with the optimizer VOQC, in
Rocq~\cite{paykin2017qwire,rand2018qwirepractice,hietala2021voqc,hietala2021proving},
CoqQ~\cite{zhou2023coqq}, Qbricks~\cite{chareton2021qbricks}, and in
Isabelle a quantum Hoare logic~\cite{liu2019quantum} and a library for
quantum information~\cite{bordg2020dirac}. VyZX~\cite{lehmann2025vyzx}
formalizes the ZX-calculus in Rocq and proves its rules sound. In this
work a rewrite rule must be sound, and soundness is a finite
computation. Completeness is the converse question, whether the rules
derive every true equation. A complete set of rules with a normal form
decides the equality of circuits, and an optimizer built on a complete
set of rules misses no equation. Our theorems are stated against
matrices, so they could be connected to the semantics of these
projects. We have not done so.

\paragraph{Completeness on paper.} The papers we formalized belong to a
larger literature, almost all of it checked by hand: complete theories
for all quantum circuits~\cite{clement2023complete} and minimal
versions of them~\cite{clement2024minimal}, multi-qutrit Clifford
circuits~\cite{li2025qutrit}, which \citet{bian2026qupit} generalize to
every odd prime dimension, and for the ZX-calculus the qubit
stabilizer, Clifford+$T$ and qupit stabilizer
fragments~\cite{backens2014zx,jeandel2018zx,booth2022qupitzx,poor2023travaganza}.
The programming-language line around the reversible language $\Pi$
reaches quantum computing by adding gates to $\Pi$, and its
completeness results for quantum fragments are translations from
presentations of matrix groups: the square-root extension
of~\citet{carette2024squareroots} builds on two presentations we
formalized, Selinger's and Bian and Selinger's, Hadamard-$\Pi$ on its
own presentation of $O_n(\Z[\frac1{\sqrt2}])$, and controlled circuits
have been axiomatized uniformly in this
setting~\cite{heunen2026onerig}. The classical core, $\Pi$ as the free
symmetric rig groupoid, was formalized in
HoTT-Agda~\cite{choudhury2022symmetries}; the quantum completeness
results were not.

\paragraph{Large formalizations of published proofs.} The four colour
theorem~\cite{gonthier2008fourcolour}, the odd order
theorem~\cite{gonthier2013oddorder} and the Kepler
conjecture~\cite{hales2017kepler} were each one large theorem,
formalized by a team over years. Our sources are the opposite: many
proofs of moderate size, which share methods and differ in their
details. A common library and generated code paid off for that reason.
The computational parts of our proofs are proofs by
reflection~\cite{boutin1997reflection}, as in the four colour theorem:
a Boolean check, proved sound once, and run by the typechecker.
\citet{ringer2019qed} survey the engineering side of such projects.

\paragraph{Language models and proof assistants.} Most work on language
models for proof assistants asks a model to prove one theorem at a time,
and measures success on a
benchmark~\cite{first2023baldur,yang2023leandojo,jiang2023dsp}. At the
other end, the formalization of Fermat's Last
Theorem~\cite{anthropic2026flt} was written by models working largely
on their own, and its statement was checked against the one in Mathlib.
Our use lies between the two. The agent worked on whole developments
for weeks, with the papers as specifications and a fixed library for
the statements, and most of its output was ordinary engineering:
definitions, scripts that generate proofs, and changes for
performance. \citet{hales2026reliability} argues that every formal
proof needs a human audit of its statement; our scaffold is a way to
make that audit short for a family of papers. We do not claim that this
experience generalizes. We report it because the time the work took
matters to anyone who plans similar work.
